% =============================================================================
% bab5-kesimpulan.tex — BAB V KESIMPULAN DAN SARAN
% Skripsi / Tesis / Disertasi (Lampiran I)
% =============================================================================

\chapter{KESIMPULAN DAN SARAN}
\label{bab:kesimpulan}

\section{Kesimpulan}
\label{sec:kesimpulan}

Berdasarkan hasil penelitian dan pembahasan yang telah diuraikan pada bab sebelumnya,
dapat ditarik beberapa kesimpulan sebagai berikut:

\begin{enumerate}
  \item Kesimpulan pertama berkaitan dengan rumusan masalah pertama. Uraikan secara
        singkat dan jelas temuan utama penelitian terkait pertanyaan penelitian pertama.

  \item Kesimpulan kedua berkaitan dengan rumusan masalah kedua. Uraikan temuan utama
        penelitian terkait pertanyaan penelitian kedua.

  \item Kesimpulan ketiga berkaitan dengan rumusan masalah ketiga. Uraikan temuan utama
        penelitian terkait pertanyaan penelitian ketiga.
\end{enumerate}

\section{Saran}
\label{sec:saran}

Berdasarkan kesimpulan di atas, peneliti memberikan beberapa saran sebagai berikut:

\subsection{Saran Teoritis}
\label{subsec:saran-teoritis}

\begin{enumerate}
  \item Bagi penelitian selanjutnya, disarankan untuk memperluas cakupan penelitian
        dengan menambahkan variabel-variabel lain yang belum diteliti dalam penelitian ini.

  \item Penelitian selanjutnya dapat menggunakan metode penelitian yang berbeda, misalnya
        pendekatan kualitatif atau metode campuran (mixed methods) untuk memperkaya
        temuan penelitian.
\end{enumerate}

\subsection{Saran Praktis}
\label{subsec:saran-praktis}

\begin{enumerate}
  \item Bagi [perusahaan/organisasi/instansi terkait], disarankan untuk \ldots

  \item Bagi pihak manajemen, disarankan untuk mempertimbangkan \ldots

  \item Bagi pemangku kepentingan lainnya, \ldots
\end{enumerate}
